\Needspace{12\baselineskip}
\SourceLeaf{93}%

\begin{center}
REMARKS ON THE »EXPANSION OF $(h-\cos\varphi)^{-\frac{3}{2}}$
INTO THE SERIES $A^{(0)}+2A^{(1)}\cos\varphi+2A^{(2)}\cos2\varphi+\cdots$«.
\end{center}

This expansion of $(h-\cos\varphi)^{-\frac{3}{2}}$ probably dates from around
1813. Its particular interest lies in its several connections with other
investigations by \textsc{Gauss}.

For the first two coefficients $A^{(0)}$, $A^{(1)}$ we begin with
the following integral representations.

\[
A^{(0)}=\int_{0}^{2\pi}\frac{d\varphi}{2\pi(h-\cos\varphi)^{\frac{3}{2}}},\qquad
A^{(1)}=\int_{0}^{2\pi}\frac{\cos\varphi\,d\varphi}{2\pi(h-\cos\varphi)^{\frac{3}{2}}}.
\]

\noindent Substituting $\varphi=2T$ and setting $h-1=m^{2}$, $h+1=n^{2}$ gives

\[
A^{(0)}=\int_{0}^{2\pi}\frac{dT}{2\pi(mm\cos T^{2}+nn\sin T^{2})^{\frac{3}{2}}},
\]

\[
A^{(1)}=\int_{0}^{2\pi}\frac{(\cos^{2}T-\sin^{2}T)\,dT}{2\pi(mm\cos T^{2}+nn\sin T^{2})^{\frac{3}{2}}}.
\]

\noindent The calculation of $A^{(0)}$ and $A^{(1)}$ from $m$ and $n$ by the
arithmetic–geometric mean algorithm was carried out by \textsc{Gauss} in Articles 16 ff.\
of his paper »Determinatio attractionis etc.« (Collected Works III p.\,331 ff.).
The integrals denoted there in Art.\,19 by $P$ and $Q$ give

\[
A^{(0)}=P+Q,\qquad A^{(1)}=P-Q.
\]

\Needspace{10\baselineskip}
\noindent To bring the notation used there into agreement with that used here in calculating the
arithmetic–geometric mean, we start
from $m^{2}=h-1$, $n^{2}=h+1$ and find in terms of $t$

\[
m=\frac{1-t}{\sqrt{2t}},\qquad n=\frac{1+t}{\sqrt{2t}}.
\]

\noindent The numbers $a$, $b$ from which the mean $a^{(\infty)}$ is formed here are
therefore

\[
a=\Bigl(\frac{1-t^{2}}{2t}\Bigr)^{2}m',\qquad b=\Bigl(\frac{1-t^{2}}{2t}\Bigr)^{2}n',
\]

\noindent and we have

\[
a^{(\infty)}=\Bigl(\frac{1-t^{2}}{2t}\Bigr)^{2}\mu,
\]

\noindent where $\mu$, as in\,the\,cited work, is the arithmetic–geometric mean of $m$ and $n$.
At the end of Art.\,17 of that paper is the equation for $\nu$,
which can be written as follows using the relations immediately preceding it there.

\[
\nu=2\frac{\lambda'}{m'}\Bigl(1+2\frac{\lambda''}{m''}\Bigl(1+2\frac{\lambda'''}{m'''}(1+\cdots\Bigr\}.
\]
